\documentclass[11pt]{article}
\usepackage[margin=1in]{geometry}
\usepackage{amsmath}
\usepackage{amssymb}
\usepackage{booktabs}
\usepackage{graphicx}
\usepackage{microtype}
\usepackage[hidelinks]{hyperref}

\newcommand{\Faq}{\textit{F.~aquilonia}}
\newcommand{\Fpo}{\textit{F.~polyctena}}
\newcommand{\Rst}{\ensuremath{R_{\mathrm{st}}}}
\newcommand{\rsq}{\ensuremath{r^{2}}}

\title{\vspace{-3em}Screening the \textit{Formica} hybrid zone for intrinsic
incompatibilities: a structure-corrected two-locus
association scan\\[0.3em]\large(work in progress)}
\author{}
\date{}

\begin{document}
\maketitle
\vspace{-3.5em}

\noindent\textit{Work-in-progress note. This document records methods, reasoning
and interim results for the intrinsic (Dobzhansky--Muller) arm of Module D, beyond
the descriptive among-population summary reported separately. Everything here is
structure-corrected but pre-simulation: a neutral, recombination-matched null
(Module E) is still required to convert the leads below into excess-over-neutral
statements. Numbers may change as that null and the analyses mature.}

\section{What an intrinsic incompatibility looks like in replicate hybrids}

The design is twenty independent, non-migrating hybrid populations, each sorting
\Faq/\Fpo{} ancestry out of the same admixture. Against that background an intrinsic
incompatibility can leave three distinct traces, and separating them is what
organises the analysis.

\paragraph{(1) Unidirectional sorting.} An incompatibility in which one parental
combination is fitter drives the \emph{same} parent's ancestry toward fixation
repeatedly across replicates. This is the per-locus parallelism that Module A
measures, and it is not a two-locus problem. Crucially it is also \emph{invisible}
to any two-locus test: linkage disequilibrium between two loci is bounded by their
allele frequencies, and as either locus approaches fixation the disequilibrium it
can carry goes to zero. The more completely a locus sorts unidirectionally, the less
LD it can express. Two-locus tests therefore cannot see the most strongly
unidirectionally-sorted loci at all; those belong to Module A, and Module D is
complementary rather than redundant.

\paragraph{(2) Bidirectional sorting (equal-fitness incompatibility).} If the two
parental combinations are \emph{equally} fit while the recombinants are not, each
replicate independently fixes one or the other compatible combination, chosen by
drift. Per locus this looks like strong sorting in \emph{random} directions --- some
populations toward \Faq, some toward \Fpo{} --- i.e.\ bidirectional sorting. Its
single-locus parallelism test is null by construction ($n_{\mathrm{aqu}}\approx
n_{\mathrm{pol}}$), so Module A cannot flag it. The signal lives at the pair level,
which is what the structure-corrected scan below is built to catch. Note that the among-population component of two-locus LD
is \emph{maximised} exactly here: a locus fixed in the same direction everywhere has
no among-population variance, whereas one fixed in different directions across
replicates has a great deal --- so bidirectional loci are the natural home of the
among-population DMI signal.

\paragraph{(3) Coupling maintained within populations.} A third trace is elevated
two-locus LD \emph{within} populations (Ohta's $D^2_{IS}>D^2_{ST}$): particular
allele combinations staying rarer than expected inside demes. This is the classic
within-deme epistatic-selection signature, but it is badly confounded by population
structure --- in particular by the fact that our samples are workers from few
colonies (nine of twenty ``populations'' are a single nest; nestmates are siblings),
so cryptic relatedness inflates apparent within-population LD (a Wahlund effect). Any
test of this trace must therefore control for relatedness.

\medskip
\noindent The descriptive among-population screen reported separately established
that the \emph{raw} among-population LD between unlinked clusters is dominated by a
single genome-wide ancestry gradient (drift/structure), with no pair specificity.
The analysis below removes that confound with a mixed model that conditions on
genome-wide relatedness (Section~\ref{sec:emmax}), addressing traces (2) and (3) at the
individual level, but must first clear a methodological obstacle: cross-chromosome
paralogy (Section~\ref{sec:paralogy}).

\section{Structure-corrected two-locus association (EMMAX)}
\label{sec:emmax}

\paragraph{Rationale and model.} The cleanest way to ask whether two loci are
associated \emph{beyond} what shared ancestry and relatedness predict is a linear
mixed model. We treat one LD cluster's consensus dosage as a quantitative trait and
test its association against every other cluster,
\begin{equation}
y_i = \mu + \beta\, x_j + u + \varepsilon, \qquad u \sim N(0, \sigma_g^2 K),\quad
\varepsilon \sim N(0, \sigma_e^2 I),
\end{equation}
where $y_i$ is the focal cluster's dosage across the 164 individuals, $x_j$ a
candidate partner, and $K$ a genomic-relationship matrix (EMMAX; Li, Kemppainen,
Rastas \& Meril\"a, a method built for LD-clustered genome-wide data). This is a
GWAS in which the ``phenotype'' is one locus and the ``markers'' are all the others.
The variance components are estimated once under the null and reused for every
partner, so a whole Manhattan is cheap.

\paragraph{Why one random effect handles both confounds.} The leading eigenvectors of
$K$ are the ancestry/structure principal components, so conditioning on $K$ removes
the genome-wide ancestry gradient --- the confound that swamped the descriptive
screen --- at the individual level. And because $K$ also encodes the sibling
relatedness among nestmates, the Wahlund inflation of within-population LD (trace~3
above) is regressed out in the \emph{same} term. A residual peak between two loci is
therefore pair-specific association with both the ancestry gradient and the colony
structure already removed. Continuous eMLG dosages are used directly, with no
hard-calling. Because the scan is symmetric we use \emph{double}
leave-one-chromosome-out: both the focal and the partner chromosome are dropped from
$K$, so neither tested locus's own region deflates its association. This recovers a
little power over single leave-one-out and leaves the test calibrated
($\lambda\approx 0.98$--$1.05$). One further point is worth stating, because it runs
against intuition: $K$ must be built from the ancestry-\emph{informative}
(differentiated) markers, not from a neutral, LD-pruned background. The confound $K$
has to absorb here is the ancestry-sorting axis, which is a \emph{selected} structure
carried by those markers; a neutral basis (right for a demographic control such as the
BayPass $\Omega$) dilutes that eigenvector and under-corrects, inflating $\lambda$ to
$\sim\!1.2$ and leaking the ancestry gradient back as spurious positive associations.
Finally, each focal trait is residualised on the top ten genome-wide PCs before the
mixed model: $K$ cannot remove a \emph{secondary} structure axis that is concentrated
on the very chromosomes double-LOCO drops, and such axes otherwise surface as hubs
(Section~\ref{sec:metanet}). Conditioning on the PCs strips the dominant, climate-aligned
axes --- after it the random-focal control gives essentially no significant partner ---
while calibration holds. A first, targeted run scanned 51 focal loci (bidirectional
candidates, extreme among-population-covariance loci, and controls); peaks were called at
a Bonferroni genome-wide threshold, and each is signed (coupling vs repulsion) by its
structure-corrected GLS coefficient --- both cis and trans hits are recovered.

\paragraph{Reading a Manhattan.} The self-peak at the focal locus and its local LD
block are the built-in positive control that the pipeline works; the signal is any
\emph{additional} peak that is unlinked from the focal locus (different chromosome,
or same chromosome $>10$~cM). The scan is well calibrated: genomic-inflation
$\lambda$ ranges 0.95--1.02 across focal sets and the background is flat
(Fig.~\ref{fig:emmax}b, a control locus with only its self-peak). Fig.~\ref{fig:emmax}a
shows a focal locus with a clean self-peak and a spread of unlinked peaks --- exactly
the structure-corrected pair-specific association we are after.

\paragraph{The first result is a warning.} The top associations are correlated
\emph{within} populations (median within-population $r\approx 0.99$), robust to
dropping the most influential individuals --- so they are neither the ancestry axis
nor a few-individual Wahlund artifact. But $r\approx 0.99$ between two unlinked loci,
sustained inside every population, is not what an incompatibility makes; it is the
signature of paralogy. This is treated next, and it is why the Manhattan in
Fig.~\ref{fig:emmax}a distinguishes clean peaks (orange) from duplicate peaks (grey
crosses): the focal locus's strongest partners are duplicates, its moderate partners
are genuine candidates.

\begin{figure}[h]
\centering
\includegraphics[width=\textwidth]{../Figures/moduleD_emmax_writeup.png}
\caption{Structure-corrected two-locus association (EMMAX, double
leave-one-chromosome-out $K$ from the differentiated eMLGs). \textbf{(a)} A focal
locus (F4250, Chr2) showing its self-peak (red) and
unlinked partner peaks, split by the paralogy filter into genuine candidates
(orange) and cross-chromosome duplicates (grey crosses). \textbf{(b)} A
consensus-bidirectional focal locus (F14324, Chr4): only its self-peak, a flat
structure-corrected background ($\lambda\approx 1$). Dashed line: Bonferroni
threshold.}
\label{fig:emmax}
\end{figure}

\section{Cross-chromosome paralogy, and why distance cannot remove it}
\label{sec:paralogy}

\paragraph{The problem.} Two loci on different chromosomes are at infinite
recombinational distance, so any strong LD between them cannot be linkage. It can
be a genuine evolutionary correlation --- but those are \emph{moderate} --- or it can
be a technical artifact in which the two ``loci'' are not distinct at all. When a
segmental duplication or a mis-assembly places one genomic region at two locations,
reads from the single true region map to both, and the two ``clusters'' become
genotype-identical by construction. Such pairs are the \emph{loudest}
inter-chromosomal disequilibrium in any dataset, and they are not biology. In the
EMMAX run they reach $\rsq\approx 0.98$ and genotype concordance $1.000$, and they
recur on particular scaffolds (Chr26). A genetic-distance rule is powerless against
them: different chromosomes already are maximal distance. (Distance does its job on
the \emph{within}-chromosome side --- our conservative complexity-reduction merge at
$0.5$~cM rather leaves nearby clusters split than over-merges them, and the
$10$~cM unlinked threshold then excludes those split neighbours from long-range
candidacy --- but that is a separate problem.)

\paragraph{The discriminant.} We flag duplicates by the median \emph{within-population}
correlation of the two clusters' dosages. Within-population removes the
among-population (ancestry-axis) component, which is not what paralogy is about;
what remains is individual-level identity. Genuine long-range coupling in these data
sits at $|r|\lesssim 0.6$; a duplication sits at $|r|\approx 1$. A flipped or
complementary duplicate (one copy allele-swapped relative to the other) gives
$r\approx -1$, which the absolute value still catches. We flag pairs with median
within-population $|r|>0.9$. Two independent signals corroborate the flag: genotype
\emph{concordance} approaching $1$ (positive duplicates), and per-cluster excess
\emph{heterozygosity} --- a collapsed duplication makes reads from two paralogs look
constitutively heterozygous, inflating $H_o$ (Fig.~\ref{fig:paralogy}b). The flag is
deliberately simple so it can be re-thresholded from the stored values; and it is a
screen, not a proof, because an extremely strong genuine within-deme incompatibility
could in principle also reach $|r|\approx 1$ --- which is why the corroborating
signals matter.

\paragraph{It affects both arms, unequally.} The filter is shared between the EMMAX
and the among-population arms. It flags $32/332$ EMMAX hits (10\%) but only
$55/401{,}945$ among-population survivors (0.01\%). The asymmetry is informative: a
duplicate \emph{is} structure-corrected individual-level LD, so it surfaces strongly
in the mixed-model arm, whereas the among-population statistic --- which asks about
allele-frequency covariance across populations --- is comparatively insensitive to
it. In the among-population arm the excess-heterozygosity signal cleanly separates
the flagged pairs ($H_o\approx 0.70$ vs $0.37$ background); in the EMMAX arm the
pre-selected focal set is already high-$H_o$, so within-population $|r|$ carries the
discrimination and $H_o$ is less informative (Fig.~\ref{fig:paralogy}). Removing the
flagged pairs leaves 300 clean EMMAX candidates. The filter is also a free
by-product: it identifies duplicated/mis-assembled regions directly from the
population data.

\begin{figure}[h]
\centering
\includegraphics[width=\textwidth]{../Figures/moduleD_paralogy_writeup.png}
\caption{The paralogy discriminant. \textbf{(a)} Distribution of within-population
$|r|$ for candidate pairs: the among-population survivors (green) concentrate at
moderate $|r|$ with a thin tail, whereas the EMMAX hits (blue) carry a heavy tail to
$|r|\approx 1$ --- the duplicates. Dashed line: the $0.9$ flag. \textbf{(b)} Pairs at
$|r|\approx 1$ (orange, flagged) also carry elevated heterozygosity, corroborating a
duplication origin; excess $H_o$ alone is not sufficient (many high-$H_o$ pairs at
moderate $|r|$ are genuine).}
\label{fig:paralogy}
\end{figure}

\section{A systematic pass: global FDR, region-merging, and the candidate modules}
\label{sec:metanet}

The first, targeted run (Section~\ref{sec:emmax}) called peaks one focal Manhattan at a
time under a per-focal Bonferroni threshold. To put every association on one footing ---
and to let the hubs and the pair-specific residue separate themselves under a single
criterion --- we re-read the whole EMMAX scan as one procedure. All $1.38\times10^{6}$ unlinked $(\text{focal},
\text{partner})$ tests were pooled and controlled at a global Benjamini--Hochberg
$q<0.01$ (which reproduces the earlier count, ${\sim}276$ pairs, but as one
whole-experiment error rate rather than $51$ independent ones). We then (i) removed
cross-locus paralogues ($|r_{\mathrm{within}}|>0.9$; $-29$ pairs); (ii) applied a
\emph{third-level merge} --- for this network only, not the canonical clustering ---
grouping same-chromosome clusters within $10$~cM by \emph{average-linkage} clustering of
their consensus correlation (cut at $|r|>0.5$), so that one low-recombination region stops
masquerading as a swarm of interacting loci. Average linkage requires the whole merged
group to be mutually correlated; connected components (the naive choice) is single linkage
and chains uncorrelated ends together, so a $16$-cluster block on Chr26 that single linkage
swallows whole instead resolves into a few mutually-correlated sub-blocks
($156\to127$ meta-nodes). We then (iii) dropped edges whose among-population
allele-frequency correlation does not survive leave-one-population-out
($|r|_{\mathrm{LOO}}<0.3$; $-50$ edges), which by itself dissolved a spurious degree-$19$
hub on Chr4 built from a single deme --- exactly the near-fixed / leverage artifact the
MAF--LD bound predicts. Low-recombination meta-nodes (recombination percentile $<0.1$)
were \emph{labelled} ``structure'' rather than removed. What remains is $110$ meta-edges
over $83$ meta-nodes ($18$ trans, $92$ cis; Fig.~\ref{fig:metanet}).

The network splits cleanly into two kinds of object. The dominant one is a single grey
\emph{structure} complex --- the low-recombination hubs
(the merged F33028/F33095 anchor and F27700, plus F57361) wired almost entirely to one
another: $66$ of the $110$ edges are structure--structure. These loci sit in the lowest
recombination percentile (centromeric and pericentromeric), which is why they persist
through PC-conditioning. What they are is better read one module at a time (the modules
below) than from the pooled set of all structure loci, which mixes several distinct
co-ancestry modules and so has no coherent joint interpretation.

\paragraph{Low recombination is a caveat, not a verdict.} We label these regions
``structure'' as an \emph{annotation}, never a filter, because low recombination does not
imply neutrality --- it is exactly where coadapted/epistatic complexes are expected to
accumulate and persist (recombination suppression is a hallmark of coadaptation). For an
unlinked pair the between-region recombination fraction is $0.5$ regardless of local rate,
so low local recombination does not ``hold'' cross-chromosome LD; it makes each region a
clean, low-noise ancestry \emph{block}, i.e.\ the highest-signal marker of whatever
residual covariation exists after PC-conditioning --- finer neutral structure \emph{or}
selection. These low-recombination regions are \emph{consistent with} neutral relic
co-ancestry but do not by themselves exclude selection (parallel or stabilising selection
also lowers differentiation). The
only clean discriminant is a recombination-matched neutral null (Module E,
Section~\ref{sec:synth}): every locus, low-recombination included, is carried into that
null and judged against its own local-rate baseline, so a low-recombination pair whose LD
\emph{exceeds} even that high baseline earns its way back as a coadaptation candidate.

That the structure ``module'' is not one block but several is confirmed directly. An
average-linkage clustering of the meta-node consensuses \emph{across} chromosomes (cut at
$|r|>0.4$; distinct from the within-chromosome merge, which only collapses redundant
neighbours) resolves the network into $23$ co-ancestry modules, of which a handful are
substantial: a dominant module of $16$ meta-nodes spanning $12$ chromosomes, and secondary
modules of $14$, $9$ and $7$ meta-nodes. Plotted one module at a time
(Fig.~\ref{fig:modhm}), each is a single coherent block --- across all its SNPs on six or
more chromosomes the individuals take one of the three dosages \emph{together}, exactly as
a co-inherited admixture haplotype block should, and unrelated to the other modules the
all-structure view had mixed in. Drawing only the within-module links on the circos (79 of
the 110 edges) makes the same point genome-wide: the co-ancestry resolves into a few
disjoint colour-coded bundles, not one hairball. The candidate F33454 aquilonia module
below falls out as its own module here too.

\begin{figure}[h]
\centering
\includegraphics[width=0.92\textwidth]{../Figures/moduleD_module_13_heatmap.png}
\caption{One cross-chromosome co-ancestry module ($9$ clusters / $7$ meta-nodes across
Chr8/10/11/13/15/21), as a genotype heatmap of its member SNPs (columns ward.D2-clustered;
colour bars: cluster, meta-node, sorting, DI; rows = individuals along the module's
co-ancestry PC1). The module behaves as a \emph{single} co-inherited haplotype block: an
individual is homozygous-reference (blue), heterozygous (pale), or homozygous-alternate
(red) consistently across all $2926$ SNPs on all six chromosomes, and the three groups do
not coincide with populations --- individual-level admixture-block co-ancestry, not
geographic structure or a pairwise incompatibility.}
\label{fig:modhm}
\end{figure}

The gold \emph{candidate} nodes are the residue of interest, and they are not merely
isolated pairs. Thirteen edges have both endpoints outside the structure module
(Table~\ref{tab:cand}), and they organise into three kinds: (a) an \emph{aquilonia
co-sorting module} anchored on F33454 (Chr10, $50$~cM, degree $12$), whose partners on
Chr3/4/5/9/11/13/19 are almost all aquilonia-sorted --- coordinated ancestry co-selection
toward the minor (aquilonia) background rather than an obvious pairwise incompatibility;
(b) a single strong, ancestry-\emph{independent} trans pair, F11431--F49480
(Chr3--Chr16, $q=2.2\times10^{-10}$, both loci unsorted) --- the cleanest pairwise-DMI
candidate the network contains; and (c) a handful of same-ancestry \emph{cis} pairs
(polyctena: F17114--F49429, F15820--F26284; aquilonia: F22294--F39587), i.e. co-sorting,
not repulsion. Stripped of structure and of ancestry co-sorting, the pair-specific
residue is thus F11431--F49480 together with a couple of weak, leverage-stable trans
pairs --- a handful of edges in all, reached by one reproducible criterion.

\begin{figure}[h]
\centering
\includegraphics[width=0.86\textwidth]{../Figures/moduleD_trans_network.png}
\caption{The Module-D meta-network after global FDR ($q<0.01$), paralogy removal,
third-level within-region merging, and the leave-one-population-out leverage filter.
Nodes are meta-nodes sized by degree; grey = low-recombination ``structure'' label,
gold = candidate region; blue edges cis (coupling), orange edges trans (repulsion). The
large grey complex (F33095/F27700/F57361) is the admixture-block structure module; the
gold hub F33454 is the aquilonia co-sorting module; scattered gold--gold pairs are the
pair-specific candidates.}
\label{fig:metanet}
\end{figure}

\begin{figure}[h]
\centering
\includegraphics[width=0.78\textwidth]{../Figures/moduleD_candidate_circos.png}
\caption{Genomic view of the candidate loci: each association drawn as a band spanning the
full extent of the two meta-nodes it joins (cis blue, trans orange; here at full opacity as
the set is sparse, with a minimum band width enforced so narrow regions stay visible). The
Chr10 anchor F33454 fans out to aquilonia-sorted partners on
Chr3/4/5/9/11/13/19 (bands landing on the blue aquilonia stretches of the sorting ring) --
the aquilonia co-sorting module -- while the single orange band F11431--F49480 (Chr3--Chr16)
is the one ancestry-independent trans pair. Outer track: $\mathrm{ld_w}>0.2$ per marker; inner
rings: genome-wide sorting, then the meta-node structure(grey)/candidate(gold) label.}
\label{fig:candcircos}
\end{figure}

\begin{figure}[h]
\centering
\includegraphics[width=0.78\textwidth]{../Figures/moduleD_structure_circos.png}
\caption{Genomic view of the structure / co-ancestry module (band opacity $\propto$ degree;
black = the F27700--F33095 heatmap pair). The module is almost entirely \emph{cis} (blue):
low-recombination admixture blocks on different chromosomes couple in the \emph{same}
ancestry direction, converging on the Chr8/Chr10 anchors. The contrast with
Fig.~\ref{fig:candcircos} -- cis-coupled co-ancestry here, the trans/repulsion signal with
the candidates -- is the visual statement that the two are different phenomena.}
\label{fig:structcircos}
\end{figure}

\begin{table}[h]
\centering
\caption{The $13$ candidate--candidate meta-edges (both endpoints outside the structure
module), with chromosomes, coupling, FDR $q$, and the ancestry-sorting class of each
endpoint. The F33454 anchor and its aquilonia-sorted partners form a co-sorting module;
F11431--F49480 is the one strong ancestry-independent trans pair.}
\label{tab:cand}
\small\setlength{\tabcolsep}{5pt}
\begin{tabular}{llccl}
\toprule
Pair & Chr & coupling & $q$ & sorting (a/b) \\
\midrule
F11431--F49480 & 3/16  & trans & $2.2\times10^{-10}$ & unsorted/unsorted \\
F10974--F33454 & 3/10  & trans & $2.7\times10^{-4}$  & aqu/aqu \\
F32802--F33454 & 9/10  & trans & $1.1\times10^{-3}$  & aqu/aqu \\
F11265--F33454 & 3/10  & trans & $6.2\times10^{-3}$  & pol/aqu \\
F15419--F33454 & 4/10  & cis   & $2.8\times10^{-11}$ & aqu/aqu \\
F33454--F57178 & 10/19 & cis   & $8.8\times10^{-5}$  & aqu/aqu \\
F33454--F41301 & 10/13 & cis   & $1.2\times10^{-4}$  & aqu/aqu \\
F33454--F35575 & 10/11 & cis   & $2.2\times10^{-4}$  & aqu/aqu \\
F17114--F49429 & 4/15  & cis   & $9.2\times10^{-4}$  & pol/pol \\
F15820--F26284 & 4/7   & cis   & $2.3\times10^{-3}$  & pol/pol \\
F21099--F33454 & 5/10  & cis   & $3.0\times10^{-3}$  & aqu/aqu \\
F22294--F39587 & 6/12  & cis   & $4.7\times10^{-3}$  & aqu/aqu \\
F45183--F45504 & 14/14 & cis   & $5.0\times10^{-3}$  & unsorted/unsorted \\
\bottomrule
\end{tabular}
\end{table}

\section{Synthesis and what the simulation will settle}
\label{sec:synth}

One structure-corrected scan, read out by a single reproducible procedure (global FDR,
average-linkage region-merging, and a leave-one-population-out leverage filter), gives a
largely negative result. Paralogy, the dominant ancestry axis, single-deme leverage, and
near-fixed loci are each removed in turn, and every removal exposes an apparent signal as
an artifact rather than an incompatibility. Two kinds of object remain. The bulk is a
low-recombination \emph{structure} module --- unlinked regions of broad co-ancestry that
resolve into several distinct co-ancestry modules rather than one interpretable block. The
residue of interest is a short candidate list: an aquilonia \emph{co-sorting} module
(F33454 and partners) that belongs with the ancestry-sorting arms A/C, a single strong
ancestry-independent trans pair (F11431--F49480), and a couple of weak, leverage-stable
trans pairs.

This is the expected place to be \emph{before} the null. Every control here removes
structure \emph{aligned} with what it conditions on; none can say whether the residue ---
or, crucially, any of the low-recombination modules --- exceeds what neutral,
recombination-matched drift would produce in twenty replicate populations of this size and
sampling design. That is Module E, and the revision that matters most is how its null is
built.

\paragraph{Low recombination must be judged, not excluded.} Low-recombination regions are
exactly where coadapted complexes are expected to accumulate, so the ``structure'' label is
an annotation, not a filter: \emph{every} locus is carried into the null. The null itself
must be matched to \emph{local} recombination --- neutral admixture simulated on the real
recombination map, with the same founder ancestries, generations and sampling --- so that a
low-recombination pair is compared against its own, legitimately high, neutral baseline and
can still surface if its LD exceeds it. Two corroborations are available even before the
simulation, and both separate selection from neutral relic co-ancestry:
\emph{replicate-population parallelism} (neutral relic LD drifts in random directions across
the twenty populations, whereas epistatic selection maintains the same combination
reproducibly) and \emph{specific exceedance} of a module over the general low-recombination
background. The coherent multi-chromosome blocks --- where an individual takes one dosage
across all of a module's SNPs on six or more chromosomes (Fig.~\ref{fig:modhm}) --- are the
first coadaptation candidates to test, not the ones to set aside. Module E converts each of
these leads into an excess-over-neutral statement, or not; until then the result is
descriptive.

\section*{Appendix: explored and set aside}

For the record, several analyses were run and are \emph{not} part of the minimal pipeline
(see \texttt{moduleD\_plan.md}). The Ohta among-population LD arm as a standalone screen: the
raw among-population signal is the ancestry cline, so it is retained only as Module E's
measurement instrument. An SNP-level bidirectional-sorting test: a weak, diffuse excess of
co-sorting whose per-locus permutation null overstates significance by not modelling shared
drift. A per-focal Bonferroni read-out: replaced by one global FDR. A neutral/LD-pruned GRM
basis: it under-corrects, because $K$ must \emph{contain} the selected ancestry axis it
removes. And a climate-association tie for the hub loci: negative --- they are among the
least climate-associated loci in the genome, so the hubs belong to neither the intrinsic nor
the extrinsic arm. These are recorded so the pruning is deliberate, not lost.

\end{document}
